% chapters/ch01.tex
\chapter{Preguntas de investigación, objetivos y justificación}

\begin{refsection}

\section{Pregunta general de investigación}

Derivada directamente del planteamiento del problema y del análisis del contexto educativo contemporáneo, la pregunta general que orienta la presente investigación doctoral es la siguiente:

\begin{quote}
¿Cómo diseñar, implementar y validar un sistema tecnológico inteligente basado en analítica de datos que permita realizar diagnósticos educativos inclusivos y apoyar la toma de decisiones pedagógicas orientadas a la equidad, la accesibilidad y la atención a la diversidad, garantizando la participación y el aprendizaje de todos los estudiantes en el contexto de la sociedad del conocimiento?
\end{quote}

Esta pregunta integra dimensiones pedagógicas, tecnológicas, metodológicas, éticas e inclusivas del fenómeno educativo contemporáneo. En particular, se sitúa en el campo de la tecnología educativa aplicada al conocimiento, donde convergen enfoques provenientes de la educación inclusiva, la analítica del aprendizaje, el Diseño Universal para el Aprendizaje (\DUA), la inteligencia artificial educativa y el desarrollo de sistemas inteligentes de apoyo a procesos pedagógicos.

La formulación de esta pregunta responde a la necesidad de explorar nuevas formas de diagnóstico educativo que superen las limitaciones de los enfoques tradicionales centrados exclusivamente en el rendimiento académico, y que permitan aprovechar el potencial de las tecnologías digitales para comprender la diversidad de los procesos de aprendizaje a partir de múltiples fuentes de información. Al mismo tiempo, reconoce que el uso de datos en educación debe estar orientado por principios de accesibilidad, privacidad, transparencia y no discriminación, de modo que la tecnología funcione como mediadora pedagógica para reducir barreras para el aprendizaje y la participación, y no como un mecanismo de exclusión.

\section{Preguntas específicas de investigación}

Con el fin de operacionalizar la pregunta general y orientar de manera precisa el desarrollo del estudio, se formulan las siguientes preguntas específicas:

\begin{itemize}
\item ¿Cuáles son los fundamentos teóricos, pedagógicos, tecnológicos y normativos que sustentan el diagnóstico educativo inclusivo y el uso responsable de datos educativos en el contexto de la sociedad del conocimiento?

\item ¿Qué variables pedagógicas, cognitivas, contextuales, socioeducativas y de accesibilidad deben considerarse para construir un diagnóstico educativo inclusivo basado en datos, evitando enfoques reduccionistas o centrados en déficits?

\item ¿Cómo diseñar un modelo conceptual y metodológico de sistema tecnológico inteligente que integre analítica de datos educativos, principios de accesibilidad y criterios de inclusión educativa, de manera que produzca diagnósticos interpretables y orientados a apoyos pedagógicos?

\item ¿De qué manera pueden integrarse técnicas de analítica del aprendizaje e inteligencia artificial en un sistema orientado al diagnóstico educativo inclusivo, asegurando criterios de explicabilidad, utilidad pedagógica y mitigación de sesgos?

\item ¿Qué principios éticos, de accesibilidad digital, privacidad y protección de datos deben incorporarse en el diseño e implementación de sistemas educativos basados en datos, considerando consentimiento informado, minimización de datos y seguridad de la información?

\item ¿Cuál es el nivel de validez, utilidad pedagógica, accesibilidad y pertinencia inclusiva del sistema tecnológico propuesto cuando se implementa en un contexto educativo real, y en qué medida apoya decisiones pedagógicas más equitativas?
\end{itemize}

Estas preguntas específicas permiten estructurar el proceso investigativo en fases coherentes que abarcan la fundamentación conceptual y normativa, la definición de variables, el diseño del modelo, el desarrollo tecnológico, la integración analítica y de inteligencia artificial, y la validación empírica del sistema en contextos reales.

\section{Objetivo general}

Diseñar, desarrollar, implementar y validar un sistema tecnológico inteligente basado en analítica de datos que permita realizar diagnósticos educativos inclusivos y apoyar la toma de decisiones pedagógicas orientadas a la equidad, la accesibilidad y la atención a la diversidad, en el contexto de la sociedad del conocimiento.

\section{Objetivos específicos}

\begin{itemize}
\item Analizar los fundamentos teóricos, pedagógicos, tecnológicos y normativos del diagnóstico educativo inclusivo y del uso responsable de datos educativos en el contexto de la sociedad del conocimiento.

\item Identificar, definir y operacionalizar variables pedagógicas, cognitivas, contextuales, socioeducativas y de accesibilidad necesarias para la construcción de diagnósticos educativos inclusivos basados en datos.

\item Diseñar un modelo conceptual y metodológico de diagnóstico educativo inclusivo que articule principios de educación inclusiva, \DUA\ y analítica del aprendizaje, orientado a la identificación de barreras, fortalezas y apoyos pedagógicos.

\item Desarrollar e implementar un prototipo funcional del sistema tecnológico propuesto, integrando técnicas de analítica del aprendizaje e inteligencia artificial, con criterios de interpretabilidad y utilidad docente.

\item Incorporar criterios de accesibilidad digital, ética de datos, privacidad y protección de la información en el diseño y funcionamiento del sistema, incluyendo mecanismos de seguridad, consentimiento y mitigación de sesgos.

\item Evaluar y validar el sistema tecnológico mediante su implementación en un contexto educativo real, analizando su validez, utilidad pedagógica, accesibilidad, pertinencia inclusiva y contribución a la toma de decisiones educativas orientadas a la equidad.
\end{itemize}

\section{Análisis del problema educativo}

El análisis del problema constituye una etapa fundamental en la formulación de la investigación, ya que permite comprender las causas, manifestaciones y consecuencias de la problemática educativa que se busca abordar. En investigación educativa, la identificación rigurosa del problema es el punto de partida para el diseño de propuestas pedagógicas y tecnológicas orientadas a la mejora de los procesos de enseñanza y aprendizaje, así como para la delimitación de un objeto de estudio viable y científicamente pertinente.

En el contexto de la educación contemporánea, caracterizado por transformaciones sociales, culturales y tecnológicas aceleradas, los sistemas educativos enfrentan el desafío de responder a una creciente diversidad de estudiantes, cuyas características, trayectorias educativas, contextos socioculturales y formas de aprender presentan una heterogeneidad significativa. Esta diversidad incluye diferencias en ritmos y estilos de aprendizaje, condiciones socioeconómicas, contextos culturales y lingüísticos, así como necesidades educativas asociadas a factores cognitivos, comunicativos, sensoriales, emocionales o contextuales. En consecuencia, el diagnóstico educativo adquiere un valor estratégico, pues permite orientar decisiones pedagógicas que reduzcan barreras para la participación y sostengan oportunidades reales de aprendizaje.

Sin embargo, en muchos contextos educativos los procesos de diagnóstico pedagógico continúan basándose predominantemente en evaluaciones estandarizadas que priorizan la medición del rendimiento académico mediante indicadores cuantitativos. Aunque estos instrumentos pueden proporcionar información útil para ciertos fines de evaluación, con frecuencia resultan insuficientes para comprender la complejidad de los procesos educativos, así como la diversidad de condiciones que inciden en el aprendizaje de los estudiantes. En particular, la literatura sobre inclusión y equidad ha señalado que tratar a todos de la misma manera puede invisibilizar desigualdades, y que la identificación de necesidades educativas debe orientarse a apoyos diferenciados, no a la categorización rígida del alumnado \parencite{unesco2020}.

En este marco, los diagnósticos educativos tradicionales suelen presentar al menos cuatro limitaciones principales. En primer lugar, tienden a centrarse en resultados académicos sin considerar de forma suficiente factores pedagógicos, contextuales y socioemocionales que influyen en el aprendizaje. En segundo lugar, suelen apoyarse en instrumentos diseñados bajo supuestos de homogeneidad del alumnado, lo cual dificulta reconocer barreras y apoyos necesarios en contextos diversos. En tercer lugar, rara vez integran de manera sistemática múltiples fuentes de información educativa, como evidencias de participación, interacción, autorregulación o condiciones de accesibilidad. En cuarto lugar, tienden a producir información poco accionable para el diseño didáctico inclusivo, lo que limita su aporte a la toma de decisiones pedagógicas.

Paralelamente, la expansión de entornos digitales de aprendizaje y plataformas educativas ha generado grandes volúmenes de datos educativos que potencialmente podrían contribuir a mejorar los procesos diagnósticos. Estos datos incluyen registros de interacción, resultados de actividades, patrones de participación, tiempos de ejecución, rutas de navegación y otras trazas digitales. No obstante, en muchos sistemas educativos estos datos no se aprovechan de manera sistemática para apoyar el diagnóstico pedagógico: la limitada integración entre pedagogía y tecnología dificulta transformar la evidencia digital en conocimiento pedagógico útil. Como resultado, los docentes pueden carecer de herramientas que les permitan interpretar información educativa compleja y convertirla en decisiones contextualizadas, sostenibles y consistentes con principios de inclusión.

Cuando los diagnósticos no logran captar la diversidad de condiciones que influyen en el aprendizaje, existe el riesgo de que las prácticas pedagógicas reproduzcan barreras para el aprendizaje y la participación. Estas barreras pueden manifestarse en estrategias didácticas poco flexibles, limitaciones para identificar necesidades educativas oportunamente, insuficiencia de apoyos o adecuaciones, y dificultades para evaluar el impacto de las intervenciones pedagógicas. De este modo, la problemática no se limita a un déficit de información, sino a la ausencia de procesos diagnósticos integrales que permitan anticipar la diversidad y orientar decisiones basadas en evidencia.

En este contexto, surge la necesidad de desarrollar enfoques diagnósticos más integrales que permitan comprender la complejidad del aprendizaje y apoyar la toma de decisiones pedagógicas informadas. La analítica del aprendizaje y la educación basada en datos constituyen campos emergentes que buscan aprovechar la evidencia generada en entornos educativos para comprender procesos de aprendizaje y apoyar decisiones pedagógicas \parencite{siemenslong2011}. Sin embargo, su aplicación plantea desafíos relacionados con la interpretación pedagógica de los datos, la protección de la privacidad, la seguridad de la información y la prevención de sesgos que puedan amplificar desigualdades.

En consecuencia, el problema central que aborda la presente investigación puede formularse como la persistencia de limitaciones en los enfoques diagnósticos tradicionales para comprender la diversidad del aprendizaje y apoyar decisiones pedagógicas inclusivas y accesibles en contextos educativos contemporáneos, a pesar de la disponibilidad creciente de datos educativos y de herramientas tecnológicas con potencial para mejorar la equidad.

La Figura \ref{fig:arbol-problema} presenta una representación sintética de esta problemática mediante un árbol del problema, en el cual se identifican causas principales que contribuyen a la situación problemática, así como algunos de sus efectos en los procesos educativos.

\begin{figure}[htbp]
\centering
\resizebox{\textwidth}{!}{%
\begin{tikzpicture}[
font=\footnotesize,
node distance=12mm and 12mm,
block/.style={draw, rounded corners=2mm, align=center, inner sep=6pt, text width=4.6cm},
arrow/.style={-{Stealth[length=2.6mm]}, thick, rounded corners=2pt}
]

\node[block] (problem)
{\textbf{Problema central}\\Diagnósticos educativos insuficientes para comprender la diversidad del aprendizaje y orientar decisiones inclusivas};

\node[block, above left=of problem] (effect1)
{Decisiones pedagógicas poco contextualizadas y de baja oportunidad};

\node[block, above=of problem] (effect3)
{Limitaciones para diseñar apoyos educativos adecuados y sostenibles};

\node[block, above right=of problem] (effect2)
{Persistencia de barreras para el aprendizaje y la participación};

\node[block, below left=of problem] (cause1)
{Uso limitado de datos educativos disponibles\\(evidencia fragmentada)};

\node[block, below=of problem] (cause3)
{Predominio de evaluaciones estandarizadas centradas en rendimiento};

\node[block, below right=of problem] (cause2)
{Escasa integración entre tecnología, pedagogía inclusiva y accesibilidad};

\node[block, below=of cause3] (cause4)
{Limitada consideración de la diversidad del alumnado y de variables contextuales};

\draw[arrow] (cause1) -- (problem);
\draw[arrow] (cause2) -- (problem);
\draw[arrow] (cause3) -- (problem);
\draw[arrow] (cause4) -- (cause3);

\draw[arrow] (problem) -- (effect1);
\draw[arrow] (problem) -- (effect2);
\draw[arrow] (problem) -- (effect3);

\end{tikzpicture}}
\caption{Árbol del problema asociado al diagnóstico educativo inclusivo.}
\label{fig:arbol-problema}
\end{figure}

\section{Justificación de la investigación}

\subsection{Justificación teórica}

La investigación se justifica teóricamente por la necesidad de articular enfoques contemporáneos de educación inclusiva, diagnóstico pedagógico y educación basada en datos en el contexto de la sociedad del conocimiento. A pesar de los avances en políticas y marcos internacionales, persisten dificultades para identificar oportunamente necesidades educativas diversas y para transformar esa identificación en apoyos pedagógicos efectivos, especialmente en contextos caracterizados por desigualdades estructurales \parencite{unesco2020}.

En este marco, el \DUA\ ofrece un enfoque para anticipar la diversidad del alumnado mediante múltiples medios de representación, acción y expresión, y compromiso \parencite{cast2018}. No obstante, su implementación requiere procesos diagnósticos que permitan reconocer barreras y necesidades de apoyo desde una lógica formativa, evitando interpretaciones rígidas o deficitarias. A su vez, la analítica del aprendizaje ha emergido como un campo relevante para comprender el aprendizaje en entornos digitales y apoyar decisiones educativas basadas en evidencia, pero enfrenta desafíos en su aplicación inclusiva, especialmente en relación con interpretabilidad y equidad \parencite{siemenslong2011}. En consecuencia, esta investigación contribuye a la construcción de un marco integrador que articula inclusión educativa, \DUA\ y analítica del aprendizaje para fundamentar diagnósticos educativos inclusivos.

\subsection{Justificación metodológica}

Desde el punto de vista metodológico, la investigación se justifica por la adopción de un enfoque mixto que integra análisis cuantitativo de datos educativos con la comprensión cualitativa de contextos, prácticas y percepciones de los actores educativos. El desarrollo del sistema tecnológico como instrumento de investigación posibilita la recolección sistemática de evidencia, la construcción de indicadores y la evaluación empírica de diagnósticos y recomendaciones en condiciones reales, coherente con investigaciones educativas complejas \parencite{creswell2018}.

\subsection{Justificación tecnológica}

La investigación se justifica tecnológicamente por la necesidad de desarrollar herramientas que transformen datos educativos en conocimiento pedagógico útil y accionable. El sistema propuesto integra analítica del aprendizaje, criterios de accesibilidad digital, ética y gobernanza de datos, y una arquitectura orientada a servicios que favorece escalabilidad, interoperabilidad y trazabilidad. De esta manera, la tecnología se concibe como un medio para apoyar procesos de diagnóstico educativo inclusivo y para sostener decisiones pedagógicas fundamentadas en evidencia, sin sustituir el juicio profesional docente.

\subsection{Justificación educativa y social}

Desde una perspectiva educativa y social, la investigación contribuye al desarrollo de herramientas que permitan reducir barreras para el aprendizaje y la participación, fortalecer la atención a la diversidad y promover prácticas pedagógicas más equitativas. Estos propósitos se alinean con el Objetivo de Desarrollo Sostenible 4, orientado a garantizar una educación inclusiva, equitativa y de calidad \parencite{onu2015}. En contextos donde la diversidad y la desigualdad coexisten, un diagnóstico inclusivo basado en datos puede apoyar la identificación oportuna de necesidades, la focalización de apoyos y el seguimiento de su efectividad, contribuyendo a disminuir brechas educativas.

\subsection{Justificación ética}

La investigación se justifica éticamente por su compromiso con el uso responsable de datos educativos y la protección de derechos de poblaciones diversas. Dado que los sistemas basados en datos pueden amplificar desigualdades si no se diseñan con criterios de justicia, transparencia y no discriminación, el estudio incorpora principios de consentimiento informado, anonimización o seudonimización, minimización de datos, seguridad de la información y prevención de sesgos algorítmicos, en coherencia con recomendaciones internacionales sobre ética en inteligencia artificial \parencite{unesco2021}.

\section{Arquitectura científica de la investigación}

Con el fin de sintetizar la estructura conceptual y metodológica de la presente investigación doctoral, se presenta a continuación un esquema que integra los principales componentes del estudio. Este esquema permite visualizar la relación entre el problema educativo identificado, el marco teórico que sustenta la investigación, el diseño metodológico adoptado y el modelo tecnológico propuesto.

La Figura \ref{fig:arquitectura-investigacion} resume la arquitectura científica de la investigación, mostrando la secuencia lógica que conecta el análisis del problema educativo con el desarrollo y validación del sistema tecnológico de diagnóstico educativo inclusivo.

\begin{figure}[htbp]
\centering
\resizebox{\textwidth}{!}{%
\begin{tikzpicture}[
font=\footnotesize,
node distance=12mm and 12mm,
block/.style={draw, rounded corners=2mm, align=center, inner sep=8pt, text width=4.9cm},
big/.style={draw, rounded corners=2mm, align=center, inner sep=10pt, text width=7.2cm},
arrow/.style={-{Stealth[length=2.6mm]}, thick, rounded corners=2pt}
]

\node[block] (problem)
{\textbf{Problema educativo}\\Diagnósticos limitados para comprender la diversidad del aprendizaje y orientar decisiones inclusivas};

\node[big, below=of problem] (theory)
{\textbf{Marco teórico}\\Inclusión educativa \quad $\cdot$ \quad \DUA\\Educación basada en datos / Analítica del aprendizaje};

\node[big, below=of theory] (method)
{\textbf{Marco metodológico}\\Enfoque mixto \quad $\cdot$ \quad \DBR};

\node[big, below=of method] (model)
{\textbf{Modelo de diagnóstico educativo inclusivo basado en datos}\\Variables multidimensionales \quad $\cdot$ \quad Recomendaciones alineadas con DUA};

\node[block, below left=of model] (system)
{Sistema tecnológico inteligente\\(arquitectura + analítica/IA + accesibilidad + ética)};

\node[block, below right=of model] (validation)
{Validación en contexto educativo real\\(utilidad, validez, accesibilidad, pertinencia)};

\node[big, below=18mm of model] (contribution)
{\textbf{Aportes científicos}\\Modelo conceptual \quad $\cdot$ \quad Sistema tecnológico\\Evidencia empírica \quad $\cdot$ \quad Conocimiento transferible};

\draw[arrow] (problem) -- (theory);
\draw[arrow] (theory) -- (method);
\draw[arrow] (method) -- (model);
\draw[arrow] (model) -- (system);
\draw[arrow] (model) -- (validation);
\draw[arrow] (system) -- (contribution);
\draw[arrow] (validation) -- (contribution);

\end{tikzpicture}}
\caption{Arquitectura científica de la investigación doctoral.}
\label{fig:arquitectura-investigacion}
\end{figure}

Como se observa en la Figura \ref{fig:arquitectura-investigacion}, la investigación se estructura a partir de la identificación de un problema educativo relacionado con las limitaciones de los diagnósticos educativos tradicionales para comprender la diversidad del aprendizaje y orientar decisiones inclusivas. A partir de este problema se construye un marco teórico que articula los enfoques de inclusión educativa, \DUA\ y educación basada en datos. Este marco conceptual sustenta el diseño metodológico de la investigación, basado en un enfoque mixto y en la \DBR. Con base en estos fundamentos se propone un modelo de diagnóstico educativo inclusivo basado en datos, implementado mediante un sistema tecnológico inteligente cuya validez se evalúa en un contexto educativo real. El resultado final es la generación de aportes científicos orientados al fortalecimiento de prácticas educativas inclusivas mediadas por tecnología.

\printbibliography[heading=subbibliography,title={Referencias (Capítulo I)}]

\end{refsection}